
% =============================================================== 6
\section{Constraints: cutting plane, penalty and barrier}

\begin{frame}{Constrained problem P4: where is the optimum?}
  \begin{columns}[T]
    \begin{column}{0.48\textwidth}
      \fig{fig_P4_feasible.png}
    \end{column}
    \begin{column}{0.50\textwidth}
      \small
      \begin{itemize}
        \item Lower cost is up-left (less area, hotter heat), but the demand constraint blocks it: the optimum sits \textbf{on the boundary} $g(x)=0$.
        \item There $\grad f$ and $\grad g$ are \textbf{parallel}:
        \[ \grad f(\xs)+\mu\,\grad g(\xs)=0,\qquad \mu\ge0 \]
        \item $\mu^\ast=5$ \textbf{k\$/yr} per unit of $g$ ($100$ m$^2$): the \textbf{shadow price} of the constraint, i.e.\ what one more unit of required area costs.
      \end{itemize}
      \begin{alertblock}{Idea shared by the next methods}
        Turn the constrained problem into easier ones: \textbf{linear programs} (cutting plane) or \textbf{unconstrained} problems (penalty, barrier).
      \end{alertblock}
    \end{column}
  \end{columns}
\end{frame}

\begin{frame}{Cutting plane (Kelley, 1960): replace $g$ by its tangents}
  \begin{columns}[T]
    \begin{column}{0.56\textwidth}
      \small
      \begin{enumerate}
        \item Solve the \textbf{LP}: minimise $c^\top x$ in the box, subject to all cuts so far (none at first).
        \item If its solution $x_k$ satisfies $g(x_k)\le\text{tol}$: \textbf{stop}.
        \item Otherwise add the \textbf{cut} (tangent plane of $g$ at $x_k$):
        \[ g(x_k)+\grad g(x_k)^\top(x-x_k)\le0 \]
        and repeat.
      \end{enumerate}
      \begin{alertblock}{Why it is safe}
        If $g$ is \textbf{convex}, its tangent lies below it: $g(x)\ge g(x_k)+\grad g(x_k)^\top(x-x_k)$. Every feasible point satisfies the cut, so cuts remove \textbf{only infeasible} points.
      \end{alertblock}
    \end{column}
    \begin{column}{0.42\textwidth}
      \small
      Here $f$ is linear, so only $g$ is linearised. For a nonlinear convex $f$: minimise a new variable $t\ge f(x)$ (the \emph{epigraph trick}) and cut $f$ as well.\\[6pt]
      \textbf{Cost:} one small LP per iteration (\texttt{scipy.optimize.linprog}).
    \end{column}
  \end{columns}
\end{frame}

\begin{frame}{Cutting plane on P4: the polygon closes in on the boundary}
  \centering
  \fig[0.98\linewidth]{fig_cutting_plane.png}\\[-1pt]
  \small Violet: region allowed by the cuts so far. Red dot: the LP solution $x_k$. Each cut slices off the previous (infeasible) vertex; iterates approach the optimum $\star$ from \textbf{outside}.
\end{frame}

\begin{frame}{Cutting plane: the code}
  \code{cutting}
  \small The Jacobian of $g$ is obtained like \texttt{numgrad}, one row per constraint (\texttt{numjac}). The LP solver is HiGHS inside \texttt{linprog}.
\end{frame}

\begin{frame}{Cutting plane: result and limits}
  \begin{columns}[T]
    \begin{column}{0.46\textwidth}
      \fig{fig_cp_conv.png}
    \end{column}
    \begin{column}{0.52\textwidth}
      \small
      \begin{itemize}
        \item 9 LPs reach $x=(2.744,\,5.917)$, i.e.\ $A=274$ m$^2$, $\Delta T=59.2$ K, violation $<10^{-4}$.
        \item Violation falls steadily; the distance to the optimum is erratic (iterates jump from vertex to vertex).
        \item The cuts accumulate: LPs grow every iteration.
      \end{itemize}
      \begin{alertblock}{Limits}
        Needs \textbf{convexity} (otherwise a cut can remove the true optimum). Slow on curved boundaries. Iterates are infeasible until the end.
      \end{alertblock}
    \end{column}
  \end{columns}
\end{frame}

\begin{frame}{Penalty and barrier: forget the constraint, change the objective}
  \centering
  \fig[0.86\linewidth]{fig_penalty_1d.png}\\
  \small Toy case: minimise $x$ subject to $x\ge1$.
  \begin{columns}[T]
    \begin{column}{0.48\textwidth}
      \begin{block}{Exterior penalty}
        $Q(x;\rho)=f(x)+\dfrac\rho2\max\{0,g(x)\}^2$\\ cost for \textbf{violating}: $\rho\uparrow\infty$.
      \end{block}
    \end{column}
    \begin{column}{0.48\textwidth}
      \begin{block}{Log barrier}
        $B(x;\mu)=f(x)-\mu\log\bigl(-g(x)\bigr)$\\ wall at the boundary: $\mu\downarrow0$.
      \end{block}
    \end{column}
  \end{columns}
\end{frame}

\begin{frame}{Exterior versus interior}
  \centering\small
  \begin{tabular}{p{3.2cm}p{5.2cm}p{5.2cm}}
    \toprule
    & \textbf{Exterior penalty} & \textbf{Log barrier}\\
    \midrule
    Iterates & \textbf{infeasible} (outside), approach from outside & \textbf{feasible} (inside), approach from inside\\
    Starting point & anything & must satisfy $g<0$\\
    Stopping early & slightly violates the constraint & a usable (feasible) design\\
    Multiplier estimate & $\hat\mu=\rho\,\max\{0,g(x_\rho)\}$ & $\hat\mu=-\mu/g(x_\mu)$\\
    Equalities & natural & not directly\\
    Price & ill-conditioned Hessian, $\sim\rho$ & ill-conditioned Hessian, $\sim1/\mu$\\
    \bottomrule
  \end{tabular}
  \vspace{0.3cm}

  \normalsize Procedure: solve $\min Q$ (or $\min B$) for a \textbf{sequence} $\rho=10,10^2,\dots$ (or $\mu=1,0.3,\dots$), each time starting from the previous solution (\emph{warm start}), with any unconstrained solver (Newton, BFGS, \dots).
\end{frame}

\begin{frame}{Penalty and barrier: the code}
  \code{penalty}
  \small \texttt{newton\_inner} is the damped Newton of the previous section (returning $+\infty$ outside the barrier's domain is how it keeps the iterates feasible).
\end{frame}

\begin{frame}{Penalty and barrier on P4}
  \centering
  \fig[0.62\linewidth]{fig_penalty_P4.png}\\[-2pt]
  \small
  \begin{tabular}{lrrrrr}
    \toprule
    exterior: $\rho$ & $10$ & $10^2$ & $10^3$ & $10^4$ & $10^5$\\
    $A$ [m$^2$] & 225 & 270 & 275 & 275 & 275\\
    $\hat\mu=\rho\,g$ & 5.000 & 5.000 & 5.000 & 5.000 & 5.000\\
    \midrule
    barrier: $1/\mu$ & $1$ & $3.3$ & $10$ & $33$ & $100$\\
    $A$ [m$^2$] & 295 & 281 & 277 & 276 & 275\\
    $\hat\mu=-\mu/g$ & 5.000 & 5.000 & 5.000 & 5.000 & 5.000\\
    \bottomrule
  \end{tabular}
  \\[2pt]\footnotesize Both estimates recover the shadow price $\mu^\ast=5$ k\$/yr. (Here $f$ and $g$ are linear in $x_1$, so $x_2$ is exact from the first subproblem.)
\end{frame}

\begin{frame}{The price: conditioning, and a warning about scaling}
  \begin{itemize}
    \item As $\rho\to\infty$ the penalty term $\frac\rho2\,\grad g\,\grad g^\top$ dominates the Hessian of $Q$: one huge eigenvalue, the rest small. The \textbf{condition number grows like $\rho$} (like $1/\mu$ for the barrier) when fewer constraints are active than variables.
    \item Consequences: Newton's linear systems become inaccurate, BFGS slows, steps get tiny near the boundary (we stopped the barrier at $\mu=0.01$ in finite-difference arithmetic).
    \item \textbf{Scale your variables} ($x_1=A/100$ m$^2$, $x_2=\Delta T/10$ K): with variables of very different size the penalty acts unevenly.
  \end{itemize}
  \begin{alertblock}{How practitioners avoid it}
    Exact penalties ($\ell_1$), the \textbf{augmented Lagrangian}, and \textbf{sequential quadratic programming} (SQP) use the multiplier estimates to reach the solution with a \emph{finite} penalty.
  \end{alertblock}
\end{frame}

\begin{frame}[fragile]{Constrained problems in \texttt{scipy}}
\begin{lstlisting}
from scipy.optimize import minimize
cons = [{"type": "ineq", "fun": lambda x: -g_solar(x)}]            # scipy wants  c(x) >= 0
res  = minimize(cost_solar, [3, 3], constraints=cons, bounds=[(0, 6), (1, 9)], method="SLSQP")
\end{lstlisting}
  \small
  \begin{center}
  \begin{tabular}{llrrr}
    \toprule
    Method & kind & $A$ [m$^2$] & $\Delta T$ [K] & work\\
    \midrule
    exact solution & -- & 275.2 & 59.37 & --\\
    cutting plane (manual) & LP sequence & 274.4 & 59.17 & 9 LPs\\
    exterior penalty, $\rho=10^5$ & unconstrained seq. & 275.2 & 59.37 & 5 Newton solves\\
    log barrier, $\mu=0.01$ & unconstrained seq. & 275.4 & 59.37 & 5 Newton solves\\
    \texttt{SLSQP} & SQP & 275.2 & 59.37 & 12 iterations\\
    \texttt{trust-constr} & interior point & 275.2 & 59.37 & 24 iterations\\
    \bottomrule
  \end{tabular}
  \end{center}
  \texttt{SLSQP} is a sequential quadratic programming method: the professional descendant of the ideas on this section.
\end{frame}
